\section{Conclusion}
\label{sec:conclusion}

This paper examines whether COVID-19 produced persistent informality effects in Peru's labor market. Using a difference-in-differences design that exploits differential occupational exposure via the \citet{dingel2020} teleworkability index, we find that more teleworkable occupations experienced 19--22 pp smaller informality increases ($p < 0.05$ in 2022 and 2024; $p < 0.10$ in all post-2020 years). This differential materialized in 2021 and showed zero attenuation through 2024 ($p = 0.937$ for the persistence test).

Three features strengthen the causal interpretation. First, the teleworkability index is pre-determined, constructed from U.S. task descriptions before the pandemic, eliminating reverse causality. Second, panel attrition is not differential by treatment status, ruling out selective survival bias. Third, the effect is robust to within-individual (fixed effects) estimation, confirming that the result is not driven by compositional changes in the workforce.

The policy implications are direct. If informality is a permanent state for workers displaced from formal employment during COVID-19, then time-limited recovery programs are insufficient. Peru's formalization policies---including SUNAFIL enforcement, tax incentives for small firms, and labor inspection---must account for the fact that a large cohort of workers may now be structurally trapped in informality. The finding that the scarring effect is concentrated in the social security dimension (rather than firm size or contract type) suggests that portable social security mechanisms, decoupled from the employer relationship, may be more effective than traditional formalization strategies.

Future research should extend this analysis to include pre-pandemic ENAHO waves (2017--2019) to formally test pre-trends, and should validate the teleworkability crosswalk using \citet{saltiel2020}'s developing-country adaptation. The interaction between informality scarring and wage dynamics---whether informal workers also experience persistent wage penalties---is an important open question.
